% Source: generalization_risk.tex
\begin{notechapter}
% ---- chapter-local definitions ----
\forgetcs\risk\newcommand{\risk}[1]{R(#1)}
\forgetcs\emprisk\newcommand{\emprisk}[1]{\hat{R}(#1)}
\forgetcs\bayesrisk\newcommand{\bayesrisk}{R^{*}}
\forgetcs\bayesestimator\newcommand{\bayesestimator}{f_{\text{Bayes}}}
\forgetcs\hypclass\newcommand{\hypclass}{\mathcal{H}}
\forgetcs\bestinclass\newcommand{\bestinclass}{f_{\hypclass}^{*}}
\forgetcs\expect\newcommand{\expect}{\mathbb{E}}
\forgetcs\indicator\newcommand{\indicator}{\mathbb{I}}
% -----------------------------------
\chapter{A Unified View of Risk}\label{ch:generalization-risk}
\chaptersubtitle{Population, empirical and generalization risk}

\begin{abstract}
    The primary goal of supervised machine learning is to build models that perform well on new, unseen data. This document presents a unified framework for quantifying and understanding this goal through the lens of \emph{risk}. We define key concepts—\textbf{Generalization Risk}, \textbf{Bayes Risk}, and the \textbf{Bayes Estimator}—and decompose the total error into interpretable components: \textbf{Approximation Risk} and \textbf{Estimation Risk}. This decomposition provides a powerful theoretical foundation for understanding fundamental trade-offs, such as bias-variance and overfitting-underfitting.
\end{abstract}

\section{The Ideal Goal: Bayes Risk and the Bayes Estimator}
\label{generalization-risk:sec:bayes}

Before considering practical constraints, we define the theoretical limit of prediction performance.

\begin{definition}[Bayes Risk]
    The \emph{Bayes Risk}, denoted \(\bayesrisk\), is the minimum possible expected prediction error achievable by \emph{any} model, given the true underlying data distribution \( P(X, Y) \). It represents the \textbf{irreducible error} inherent in the problem, often due to noise in the data.
    \[
    \bayesrisk = \inf_{f} \expect_{(X,Y)\sim P(X,Y)}[L(f(X), Y)]
    \]
\end{definition}

\begin{definition}[Bayes Estimator]
    The \emph{Bayes Estimator}, denoted \(\bayesestimator\), is the ideal decision function that achieves the Bayes Risk. It is the function we would use if we had perfect knowledge of \( P(X, Y) \).
    \[
    \bayesestimator = \arg\min_{f} \expect_{(X,Y)\sim P(X,Y)}[L(f(X), Y)]
    \]
\end{definition}

The form of the Bayes Estimator depends on the choice of loss function \( L \).

\begin{example}[Squared Error Loss for Regression]
    For regression tasks with squared error loss \( L(f(X),Y)=(f(X)-Y)^{2} \), the Bayes Estimator is the conditional expectation:
    \[
    \bayesestimator(X) = \expect[Y \mid X]
    \]
    This minimizes the expected squared error because, for any \( X \), the conditional risk is:
    \[
    \expect[(f(X)-Y)^{2} \mid X] = \text{Var}(Y \mid X) + (f(X)-\expect[Y \mid X])^{2}
    \]
    which is minimized when \( f(X) = \expect[Y \mid X] \).
\end{example}

\begin{example}[0-1 Loss for Classification]
    For classification tasks with 0-1 loss \( L(f(X),Y)=\indicator(f(X)\neq Y) \), the Bayes Estimator is the mode of the posterior distribution (the maximum a posteriori classifier):
    \[
    \bayesestimator(X) = \arg\max_{y \in \mathcal{Y}} P(Y=y \mid X)
    \]
    This strategy minimizes the probability of misclassification. The derivation follows from minimizing the conditional expected loss:
    \[
    \expect[L(f(X),Y) \mid X] = \expect[\indicator(f(X)\neq Y) \mid X] = 1 - P(Y=f(X) \mid X)
    \]
    To minimize this, we maximize \( P(Y=f(X) \mid X) \).
\end{example}

In practice, \( P(X, Y) \) is unknown, so we cannot directly use the Bayes Estimator. Instead, we must \emph{learn} a model from a finite sample of data.

\section{The Practical Approach: Empirical Risk Minimization (ERM)}
\label{generalization-risk:sec:erm}

Since we cannot compute the true risk, we use the available data to approximate it.

\begin{definition}[Population Risk (Generalization Risk)]
    The true risk of a model \( f \) is its expected loss on the data distribution:
    \[
    \risk{f} = \expect_{(X,Y)\sim P(X,Y)}[L(f(X), Y)]
    \]
\end{definition}

\begin{definition}[Empirical Risk (Training Risk)]
    A Monte Carlo estimate of the population risk using a finite dataset \( d = \{(x_i, y_i)\}_{i=1}^n \):
    \[
    \emprisk{f} = \frac{1}{n} \sum_{i=1}^n L(f(x_i), y_i)
    \]
\end{definition}

The \textbf{Empirical Risk Minimization (ERM)} principle selects the model that minimizes this empirical risk from a predefined set of candidate models, the \textbf{hypothesis class} \( \hypclass \):
\[
\hat{f} = \arg\min_{f \in \hypclass} \emprisk{f}
\]

The restriction to \( \hypclass \) is necessary to make the learning problem tractable and to avoid overfitting. It encodes our prior beliefs about the structure of the solution.

\section{Generalization Risk: Model vs. Algorithm}
\label{generalization-risk:sec:model_vs_alg}

We must distinguish between the performance of a single, fixed model and the average performance of a learning procedure.

\subsection{Risk for a Specific Model}
Consider a \textbf{fixed training set} \( d \). A learning algorithm\footnote{
    The learning algorithm is the mapping 
    \( \mathcal{A}:\mathcal{D}\rightarrow\mathcal{H} \)
    that takes datasets to models.}
    \( \mathcal{A} \) produces a 
    \textbf{specific model} \( f_d = \mathcal{A}(d) \). 
    The generalization risk for this model is a deterministic quantity:
\[
\risk{f_d} = \expect_{(X,Y)\sim P(X,Y)}[L(f_D(X), Y)|D=d] =
\expect_{(X,Y)\sim P(X,Y)}[L(f_d(X), Y)]
\]
This measures the expected performance of this single model on new data.

\subsection{Expected Risk for a Learning Algorithm}
Now, consider the \textbf{training set \( D \)} as a random variable, drawn i.i.d. from \( P(X, Y)^n \). The resulting model \( f_D \) is also a random variable. The \textbf{expected generalization risk} averages over all possible training sets:
\[
\expect_{D}[\risk{f_D}] 
= \expect_{D}[\expect_{(X,Y)}[L(f_D(X), Y)|D]]
= \expect_{D, (X,Y)}[L(f_D(X), Y)]
\]
This measures the average performance of the \textbf{learning algorithm} \( \mathcal{A} \), not just of one particular model. It is the quantity we aim to minimize when designing or selecting a learning algorithm.
\section{Decomposition of Generalization Risk}
\label{generalization-risk:sec:decomposition}

Why does the model \( f_d \) learned from a finite dataset not achieve the Bayes Risk? We can decompose its error into three fundamental components.

Let \( \bestinclass \) be the best possible model within our hypothesis class (assuming the minimum is attained; otherwise replace \( \risk{\bestinclass} \) by \( \inf_{f\in\hypclass}\risk{f} \) below):
\[
\bestinclass = \arg\min_{f \in \hypclass} \risk{f}
\]

The generalization risk of a specific model \( f_d \) decomposes as follows:
\begin{equation}
\risk{f_d} = \underbrace{\risk{f_d} - \risk{\bestinclass}}_{\text{Estimation Risk}} + \underbrace{\risk{\bestinclass} - \bayesrisk}_{\text{Approximation Risk}} + \underbrace{\bayesrisk}_{\text{Bayes Risk}}
\label{generalization-risk:eq:decomposition}
\end{equation}

\subsection{Components of the Decomposition}

\begin{description}
    \item[Bayes Risk (\( \bayesrisk \))] The irreducible error, as defined in Section~\ref{generalization-risk:sec:bayes}. It is the absolute lower bound on the error for any learning algorithm.

    \item[Approximation Risk (\( \risk{\bestinclass} - \bayesrisk \))] The error incurred by restricting ourselves to the hypothesis class \( \hypclass \). It measures how poorly even the best model in \( \hypclass \) can approximate the Bayes optimal predictor \( \bayesestimator \). A rich, complex \( \hypclass \) (e.g., a large neural network) has low approximation risk.

    \item[Estimation Risk (\( \risk{f_d} - \risk{\bestinclass} \))] The error incurred because we have only a finite sample to learn from. Instead of finding \( \bestinclass \), the ERM principle finds \( f_d \), which is an estimate of it based on the empirical data. This component is related to overfitting. A simple, restricted \( \hypclass \) (e.g., linear models) typically leads to lower estimation risk.
\end{description}

\section{Estimation and Trade-offs}
\label{generalization-risk:sec:estimation_tradeoffs}

\subsection{Estimating Generalization Risk}
In practice, \( \risk{f_d} \) is unknown but can be estimated using a held-out \textbf{test set} \( \{(x_i, y_i)\}_{i=1}^m \), drawn i.i.d. from \( P(X, Y) \) \emph{independently of the training set} \( d \):
\[
\hat{R}_{\text{test}}(f_d) = \frac{1}{m} \sum_{i=1}^m L(f_d(x_i), y_i)
\]
Conditionally on \( d \), this is an unbiased and consistent Monte Carlo estimate of \( \risk{f_d} \). The training error \( \emprisk{f_d} \), by contrast, is in general optimistically biased, because \( f_d \) was chosen to make it small. Likewise, once the test set has been used to select among models, it no longer gives an unbiased estimate for the selected model. The quality of the estimate improves as the test set size \( m \) increases.

\subsection{The Fundamental Trade-off}
The choice of hypothesis class \( \hypclass \) governs a critical trade-off, directly visible in the decomposition \eqref{generalization-risk:eq:decomposition}:
\begin{itemize}
    \item A \textbf{large/rich} \( \hypclass \) reduces \textbf{Approximation Risk} but increases \textbf{Estimation Risk} (leading to potential \emph{overfitting}).
    \item A \textbf{small/simple} \( \hypclass \) reduces \textbf{Estimation Risk} but increases \textbf{Approximation Risk} (leading to potential \emph{underfitting}).
\end{itemize}
The art of machine learning lies in navigating this trade-off, for example, via model selection and regularization.

\section{Conclusion}
\label{generalization-risk:sec:conclusion}

This document has established a rigorous framework for understanding risk in machine learning:
\begin{enumerate}
    \item The \textbf{Bayes Risk} and \textbf{Bayes Estimator} define the theoretical limit of prediction.
    \item The \textbf{Empirical Risk Minimization} principle provides a practical learning rule, connecting the theoretical ideal to practical algorithms.
    \item The \textbf{Generalization Risk} of a resulting model can be cleanly decomposed into \textbf{Estimation Risk}, \textbf{Approximation Risk}, and the irreducible \textbf{Bayes Risk}.
    \item A critical distinction exists between the risk of a single model \( \risk{f_d} \) and the expected risk of a learning algorithm \( \expect_D[\risk{f_D}] \).
\end{enumerate}

This framework not only explains \emph{why} models make errors but also provides the conceptual tools to analyze and improve machine learning algorithms by explicitly considering the trade-offs between data, model complexity, and performance.
\end{notechapter}
